\documentclass[11pt]{amsart}

\usepackage[utf8]{inputenc}

\usepackage{amssymb,amsmath,mathtools}
\usepackage{booktabs,longtable,array}
\usepackage[hidelinks]{hyperref}


% ---------------------------------------------------------------------------
% Draft status (PLAN.md, track P). Authorship, venue and the acknowledgement
% of tool use are decisions for the author (PLAN.md 4.4). Every number in
% this text is produced by a test of the package or by paper/reproduce.py.
% ---------------------------------------------------------------------------

\newtheorem{theorem}{Theorem}[section]
\newtheorem{proposition}[theorem]{Proposition}
\newtheorem{lemma}[theorem]{Lemma}
\newtheorem{corollary}[theorem]{Corollary}
\theoremstyle{definition}
\newtheorem{definition}[theorem]{Definition}
\newtheorem{computation}[theorem]{Computation}
\newtheorem{observation}[theorem]{Observation}
\newtheorem{conjecture}[theorem]{Conjecture}
\newtheorem{question}[theorem]{Question}
\theoremstyle{remark}
\newtheorem{remark}[theorem]{Remark}
\newtheorem{example}[theorem]{Example}

\setlength{\emergencystretch}{3em}
\newcommand{\bbcells}{\texttt{bbcells}}
\newcommand{\code}[1]{\texttt{#1}}
\newcommand{\PP}{\mathbb{P}}
\newcommand{\RP}{\mathbb{RP}}
\newcommand{\QQ}{\mathbb{Q}}
\newcommand{\ZZ}{\mathbb{Z}}
\newcommand{\RR}{\mathbb{R}}
\newcommand{\LL}{\mathbb{L}}
\newcommand{\GW}{\mathrm{GW}}
\newcommand{\pr}{\mathrm{pr}}
\newcommand{\wt}{\mathrm{wt}}
\DeclareMathOperator{\supp}{supp}
\DeclareMathOperator{\Stab}{Stab}
\DeclareMathOperator{\Ann}{Ann}
\DeclareMathOperator{\Gr}{Gr}
\DeclareMathOperator{\Bl}{Bl}
\DeclareMathOperator{\codim}{codim}

\title[Bia{\l}ynicki-Birula cells of spherical varieties]{Bia{\l}ynicki-Birula
cells of spherical varieties from combinatorial data:\\ a companion to the
\bbcells\ package}

% Authorship to be decided (PLAN.md 4.4).
\author{Konrad Voelkel}
\date{Draft, \today}

\begin{document}

\begin{abstract}
We describe \bbcells, a pure-Python package that computes
Bia{\l}ynicki-Birula decompositions of smooth projective varieties with a
torus action from combinatorial data (fans, root data, Satake diagrams,
spherical data), together with the invariants that only depend on the
cells: Chow motives, Hodge numbers, classes in $K_0(\mathrm{Var})$, point
counts and quadratic Euler characteristics. Every front end is tested
against an independent oracle and against the holomorphic Lefschetz
formula.

The main new points are the following.
(1) Fixed-point data of smooth complete spherical varieties are assembled
orbit by orbit; we prove a formula for the normal weights at the base
points of the orbits, and derive all satellites of complete symmetric
varieties from their Satake diagrams, for every real form including
$E_6$, $E_7$, $E_8$.
(2) Where the formula leaves an unknown, a certificate based on point
counts of orbit closures decides it; all Hermitian cases we tested are
certified.
(3) Cell counts of complete symmetric varieties of type $E_7$ and $E_8$
are computed without listing fixed points, up to $2.8\cdot 10^7$ fixed
points.
(4) Beyond GKM theory, Brion's description of equivariant cohomology yields
characteristic numbers (Chasles' $3264$, Schubert's $666841088$, and a value
for quadrics in $\PP^5$ that differs from the published one), shows that
complete conics admit no canonical (Goldin--Tolman) classes for any generic
cocharacter, and that the divisors do not generate the cohomology of
complete quadrics in $\PP^4$.
(5) For the real points we compute the cellular chain complex of real toric
varieties exactly by discrete Morse theory, confirm a GKM form of
Kocherlakota's incidence rule entry by entry, and extend it to graded
decompositions of non-GKM varieties.
\end{abstract}

\maketitle
\tableofcontents

% ===========================================================================
\section{Introduction}\label{sec:intro}

Let $X$ be a smooth projective variety over a field $k$ with an action of
a split torus $T$ with finitely many fixed points. For a generic
cocharacter $\lambda$, the Bia{\l}ynicki-Birula theorem \cite{BB73}
decomposes $X$ into locally closed affine spaces, the plus-cells
$C_p = \{x : \lim_{t\to 0}\lambda(t)x = p\}$, one for each fixed point $p$,
of dimension the number of $\lambda$-positive tangent weights at $p$. The
cell counts $c_d = \#\{p : \dim C_p = d\}$ determine the Chow motive
\cite{Brosnan}, the Hodge numbers, the class in $K_0(\mathrm{Var}_k)$ and
the point counts, and they are the first layer of the motivic cell
structures studied in \cite{DuggerIsaksen,Wendt,VoelkelQuadrics,VoelkelSpherical}.

All of this depends only on combinatorial data: the fixed points and the
multisets of tangent weights (\emph{fixed-point data}). For toric and flag
varieties these data are classical. For spherical varieties in general they
are not: fixed points lie in non-closed orbits, and the tangent weights at
such points involve the normal directions of the orbits. The package
\bbcells{} computes fixed-point data for a range of spherical varieties and
everything that can be read off from them. This paper documents the
methods and the results that are new, as far as we know; the code, the
tests and three technical notes (\code{docs/spherical.md},
\code{docs/cohomology.md}, \code{docs/real.md}) contain the details.

\subsection*{Results}
\begin{enumerate}
\item \emph{Spherical varieties orbit by orbit} (Section~\ref{sec:spherical}).
The fixed points of a $G$-orbit $G/H$ are $W/W_H$ when $H$ contains a
maximal torus (Lemma~\ref{lem:orbit}). For a wonderful variety the
remaining unknowns are the normal weights at the base points of the
orbits. Proposition~\ref{prop:normal} determines them up to the
$W_H$-invariants $V$ of the Levi part of the satellite; under condition (R)
($V=0$, equivalently: the satellite has as many colours as its rank, by
Knop) they are the $W_L$-averages $\pr(-\gamma)$. For complete symmetric
varieties we derive all satellites from the Satake diagram, for every real
form (Section~\ref{subsec:symmetric}).
\item \emph{Certificates} (Section~\ref{sec:certificates}). Beyond (R), the
unknown normal weight is $\pr(-\gamma)+c\zeta$ with an integer $c$. The BB
counts of an orbit closure are a step function of $c$ and must equal its
E-polynomial, which is known independently (Brion--Peyre and
inclusion--exclusion, or Poincar\'e duality). This certifies $c=0$ in
all Hermitian cases we tested, including those where the
Atiyah--Bott--Berline--Vergne identities degenerate.
\item \emph{Large Euler characteristics} (Section~\ref{sec:large}).
Streaming over cosets gives the cell counts of the complete symmetric
varieties EI, EII, EV, EVI, EVII and EIX, up to $28\,373\,976$ fixed points.
\item \emph{Equivariant cohomology beyond GKM} (Section~\ref{sec:cohomology}).
From fixed-point data alone we reconstruct the components of $X^{\ker\chi}$
that enter Brion's description of $H_T^*(X)$ \cite{Brion98}, compute
equivariant classes of line bundles and characteristic numbers of complete
quadrics, and study cell classes and the subalgebra generated by divisors.
\item \emph{Real points} (Section~\ref{sec:real}). We compute the real cell
complex of toric varieties exactly by discrete Morse theory and use it,
together with known real loci of non-GKM varieties, to test an incidence
rule predicted by the $\eta$-attaching maps of motivic cell structures.
\end{enumerate}

\subsection*{Conventions}
Characters and cocharacters are integer vectors paired by the dot product;
for root data the characters are written in the basis of simple roots
(the torus of the adjoint group). A statement called \emph{theorem},
\emph{proposition} or \emph{lemma} is proved here or in the cited source. A
\emph{computation} reports the output of the package, reproducible with
\code{paper/reproduce.py}; an \emph{observation} is a pattern found in
computations; a \emph{conjecture} is stated as such.

% ===========================================================================
\section{Fixed-point data and cell invariants}\label{sec:basics}

\begin{definition}
\emph{Fixed-point data} of rank $r$ and dimension $n$ consist of a finite
set $F$ and, for each $p\in F$, a multiset $\wt(p)$ of $n$ nonzero vectors in
$\ZZ^r$, the tangent weights. A cocharacter $\lambda\in\ZZ^r$ is
\emph{generic} if $\langle\lambda,w\rangle\ne 0$ for all weights. The
\emph{cell dimension} of $p$ is $d_\lambda(p) = \#\{w\in\wt(p) :
\langle\lambda,w\rangle>0\}$, and the \emph{cell counts} are
$c_d = \#\{p : d_\lambda(p)=d\}$.
\end{definition}

For the fixed-point data of a smooth projective $T$-variety with isolated
fixed points, the plus-cells form a filtrable decomposition into affine
spaces \cite{BB73}, and hence
\begin{gather*}
[X] = \sum_d c_d\,\LL^d \in K_0(\mathrm{Var}_k),\qquad
M(X) = \bigoplus_d \ZZ(d)[2d]^{\oplus c_d},\\
h^{p,q}(X) = \delta_{pq}\,c_p ,
\end{gather*}
the point count over $\mathbb F_q$ is $\sum_d c_d q^d$, and the quadratic
Euler characteristic is $\chi^{\mathbb A^1}(X) = \sum_d c_d\langle -1\rangle^d
\in \GW(k)$ \cite{Levine,Hoyois}, whose signature is $\chi(X(\RR))$.

\subsection{Consistency checks}\label{subsec:checks}
Wrong fixed-point data often still produce plausible numbers, so every
computation is checked in three ways.
\begin{enumerate}
\item The cell counts do not depend on the generic cocharacter.
\item They are palindromic (Poincar\'e duality).
\item The holomorphic Lefschetz formula \cite{AtiyahBottLefschetz}: the
rational function
\[
\sum_{p\in F}\ \prod_{w\in\wt(p)} \frac{1+y\,e^{-w}}{1-e^{-w}}
\]
on $T$ must be the constant $\chi_y(X) = \sum_d c_d(-y)^d$. It is evaluated
modulo a large prime at pseudo-random points of $T$.
\end{enumerate}
The third check is much stronger than the first two for data assembled
orbit by orbit: flipping the signs of the normal weights of both
intermediate orbits of complete conics keeps (1) and (2) intact but fails
(3).

\subsection{Front ends and oracles}\label{subsec:frontends}
Table~\ref{tab:frontends} lists the front ends that produce fixed-point
data, and the independent oracle each one is tested against.

\begin{table}[ht]
\caption{Front ends and their oracles.}\label{tab:frontends}
\begin{tabular}{p{0.38\textwidth}p{0.55\textwidth}}
\toprule
front end & oracle \\
\midrule
smooth complete toric varieties (fans) & $h$-vector of the fan \\
$G/P$, all Cartan types & $\prod[d_i]_q/\prod[d_i^L]_q$ from the degrees of $W$ and $W_L$ \\
quadrics $Q_n$ & the original script \code{quadric.py}; $\sum q^i$ (+$q^{n/2}$) \\
wonderful compactifications of $G_{\mathrm{ad}}$ & orbit decomposition, $|X(\mathbb F_q)|=\sum_I |G/P_I|^2 q^{N_I}(q-1)^{|I|}\sum_{W_I}q^{\ell(w)}$ \\
toroidal horospherical varieties & Batyrev--Moreau \cite{BatyrevMoreau} \\
rank-one two-orbit completions (Knop's table) & $[X]=[\overline X]-[D]$, e.g.\ $[\mathbb{OP}^2]=\LL^8+\LL^{12}+\LL^{16}$ \\
blow-ups along $T$-stable centres & blow-up formula; complete conics $=\Bl_{v_2(\PP^2)}\PP^5$ \\
wonderful and symmetric varieties & Brion--Peyre orbit counts (Section~\ref{subsec:orbitcounts}) \\
\bottomrule
\end{tabular}
\end{table}

Two technical points recur. First, $G/P$ is never enumerated through all of
$W$: the orbit $W\cdot\omega_I$ is found by breadth-first search. Second, a
subtorus $T'\subset T$ has the same fixed points as $T$ as long as no
tangent weight restricts to zero; this is how rank-one completions
$\overline X = G'/P'$ are restricted from a larger group.

% ===========================================================================
\section{Spherical varieties orbit by orbit}\label{sec:spherical}

\subsection{Fixed points of orbits}
\begin{lemma}\label{lem:orbit}
Let $O = G/H$ be an orbit of a connected reductive group $G$ with maximal
torus $T$.
\begin{enumerate}
\item $O^T\ne\emptyset$ if and only if $H$ contains a maximal torus of $G$.
\item If $T\subset H$, then $O^T = N_G(T)H/H\cong W/W_H$ with
$W_H = N_H(T)/T$.
\item If $T\subset H$, the tangent weights at $wH$ are
$w(\Phi\smallsetminus\Phi_H)$.
\end{enumerate}
\end{lemma}
\begin{proof}
If $gH\in O^T$, then $g^{-1}Tg$ and $T$ are maximal tori of $H$, so
$g^{-1}Tg = hTh^{-1}$ for some $h\in H$ and $gh\in N_G(T)$. This gives (1)
and (2). The weights of $\mathfrak g/\mathfrak h$ are
$\Phi\smallsetminus\Phi_H$, and $0$ does not occur since
$\mathfrak t\subset\mathfrak h$.
\end{proof}

For a smooth complete $G$-variety $X$ it follows that
\[
X^T = \bigsqcup_{O :\, O^T\ne\emptyset} W/W_{H_O},\qquad
\wt(w x_O) = w(\Phi\smallsetminus\Phi_{H_O}) \cup w(N_O),
\]
where $x_O$ is a base point of $O$ fixed by $T$ and $N_O$ is the multiset of
$T$-weights of the normal space to $O$ at $x_O$, which is $W_{H_O}$-invariant.
So the fixed-point data are determined by one datum
$(\Phi_{H_O}, W_{H_O}, N_O)$ per orbit with fixed points.

\subsection{Cosets without enumerating the Weyl group}
The group $W_H$ is generated by the reflections in the roots of $H$ and
possibly by further elements; for $\mathrm{AIII}(p,p)$ it contains the
swap of the two blocks. For a reflection subgroup $W'\subset W$, Dyer
\cite{Dyer} shows that every coset $wW'$ contains a unique element of
minimal length, characterized by $w(\beta)>0$ for all positive roots
$\beta$ of $W'$. The following version of Deodhar's lemma \cite{Deodhar}
holds for reflection subgroups by the same criterion: if $w$ is minimal in
$wW'$, then either $s_iw$ is minimal in $s_iwW'$, or $w^{-1}\alpha_i$ is a
positive root of $W'$ and $s_iwW' = wW'$. A breadth-first search by length
over the inverse permutations of the roots therefore produces all minimal
representatives together with reduced words, without listing $W$. Further
generators of $W_H$ then glue these cosets. This is what makes $E_7$ and
$E_8$ accessible (Section~\ref{sec:large}).

\subsection{Wonderful varieties}\label{subsec:wonderful}
Let $X$ be a wonderful $G$-variety with spherical roots $\Sigma$, and let
$S^p$ be the set of simple roots of the parabolic subgroup $P(X)$ (the
stabilizer of the open $B$-orbit) \cite{Luna01,Timashev,Perrin}. The $G$-orbits of $X$
are the $O_I$, $I\subseteq\Sigma$, where $O_I$ is the open orbit of the
wonderful subvariety with spherical roots $I$. With
$S_I = S^p\cup\supp I$, the orbit $O_I$ is parabolically induced,
\[
O_I = G\times^{P_{S_I}} \bigl(L_{S_I}/H_{L,I}\bigr),
\]
where the \emph{satellite} $L_{S_I}/H_{L,I}$ is a cuspidal spherical
homogeneous space with spherical roots $I$. By Lemma~\ref{lem:orbit}, $O_I$
has fixed points if and only if $H_{L,I}$ has full rank in $L_{S_I}$. The
roots of the stabilizer of the base point $x_I$ in the fibre over
$eP_{S_I}$ are then $(\Phi^+\smallsetminus\Phi_{L_{S_I}})\cup\Phi_{H_{L,I}}$,
and $W_{H} = W_{H_{L,I}}$. The normal space at $x_I$ is the sum of the normal
lines $\chi_\gamma$ of the boundary divisors $D_\gamma$,
$\gamma\in\Sigma\smallsetminus I$.

\begin{proposition}\label{prop:normal}
Let $Z = Z(L_{S_I})^\circ$, let $\pr$ be the orthogonal projection onto
$\Phi_{L_{S_I}}^{\perp}$ (the average over $W_{L_{S_I}}$), and let $V$ be the
space of $W_H$-invariants in $\operatorname{span}_\QQ\Phi_{L_{S_I}}$. Then
for every $\gamma\in\Sigma\smallsetminus I$:
\begin{enumerate}
\item $\chi_\gamma|_Z = -\gamma|_Z$;
\item $\chi_\gamma$ is $W_H$-invariant;
\item hence $\chi_\gamma\in\pr(-\gamma)+V$. In particular, if $V=0$
(condition \textup{(R)}), then $\chi_\gamma = \pr(-\gamma)$.
\end{enumerate}
\end{proposition}
\begin{proof}
(2) The group $N_H(T)$ acts on the normal line of $D_\gamma$ at $x_I$.

(1) The fibre $F = L_{S_I}\cdot x_I$ over $eP_{S_I}$ is fixed pointwise by
$Z$, because $Z$ is central in $L_{S_I}$ and contained in
$T\subset H$. Its closure $\overline F$ is stable under $P_{S_I}$, since the
unipotent radical of $P_{S_I}$ acts trivially on $F$. Hence $\overline F$
contains a $B$-fixed point, and the only $B$-fixed point of a wonderful
variety is the base point $z$ of the closed orbit. So $x_I$ and $z$ lie in
one connected component of $X^Z$. Along this component, $Z$ acts on the
fibres of the $G$-linearized line bundle $\mathcal O(D_\gamma)$ by a single
character. At a point of $D_\gamma$ the fibre is the normal line of
$D_\gamma$; at $z$ the character is $-\gamma|_Z$, the $T$-weight of the
coordinate of the slice through $z$ that cuts out $D_\gamma$ (this
orientation agrees with \cite[\S1.3]{BrionJoshua}, where the weight is
$+\gamma$ at the $B^-$-fixed point).

(3) By (2), $\chi_\gamma - \pr(\chi_\gamma)$ lies in $V$, and by (1),
$\pr(\chi_\gamma) = \pr(-\gamma)$, since $\pr$ is determined by the
restriction to $Z$.
\end{proof}

By Knop \cite[Prop.~2.1]{KnopRoots}, $\dim V$ is the number of colours of
the satellite minus its rank; so (R) says that the satellite has exactly
$|I|$ colours.

\begin{remark}
Beyond (R) the $W_L$-average is wrong in general. For a satellite $GL_2/T$
with root $\alpha$, the two fixed points have normal weights
$-\gamma+a\alpha$ and $-\gamma+(\langle\gamma,\alpha^\vee\rangle-a)\alpha$,
where $a$ is the Cartan pairing of one of the two colours with $\gamma$,
and these differ in general. The package refuses such orbits unless the
remaining unknown is certified (Section~\ref{sec:certificates}).
\end{remark}

\begin{example}[Complete quadrics]
For complete quadrics in $\PP^{n-1}$ ($PGL_n/PO_n$), $\Sigma = \{2\alpha_i\}$
and $S^p=\emptyset$. The orbit of type $I$ has $L_{S_I}=S(\prod GL_{k_j})$
over the blocks of $I$ and $H_{L,I}=S(\prod O_{k_j})$, which has full rank
if and only if all $k_j\le 2$. So $\chi(X)=\sum n!/2^{\#\{j:k_j=2\}}$ over the
compositions of $n$ into parts $1$ and $2$: $3, 12, 66, 450, 3690$ for
$n=2,\dots,6$; these are the ``barred permutations'' of
\cite{BanerjeeCanJoyce}. The fixed-point data agree exactly, as multisets of
weight multisets, with independent constructions: complete conics as a
blow-up of $\PP^5$, and complete quadrics in $\PP^3$ as Vainsencher's two
blow-ups \cite{Vainsencher}.
\end{example}

\subsection{Complete symmetric varieties}\label{subsec:symmetric}
Let $\theta$ be an involution of $G$, given by its Satake diagram (black
nodes $S_\bullet$, arrows) \cite{Araki}. The complete symmetric variety of
De Concini--Procesi \cite{DCP83} is the wonderful compactification of
$G_{\mathrm{ad}}/G_{\mathrm{ad}}^\theta$. The package derives all its
fixed-point data from the diagram.
\begin{enumerate}
\item On characters $\theta = -w_\bullet\circ\varepsilon$, where $\varepsilon$
follows the arrows on white nodes and acts on black nodes by the opposition
involution of $S_\bullet$, and $w_\bullet$ is the longest element of
$W_{S_\bullet}$. The spherical roots are the $\alpha-\theta(\alpha)$, one for
each arrow orbit of white simple roots, and $S^p = S_\bullet$.
\item The satellite of $I$ is the symmetric pair on the subdiagram
$S_\bullet\cup\supp I$. It has fixed points if and only if every component
of this subdiagram is $\theta$-stable of equal rank, i.e.\ $\varepsilon=-w_0$
on it.
\item An equal-rank involution is $\mathrm{Ad}(t_j)$ with $\alpha_j(t_j)=-1$
for a Borel--de Siebenthal node $j$ \cite{BorelDeSiebenthal} and
$\alpha_i(t_j)=1$ otherwise. The node is found by recognizing the
subdiagram up to diagram automorphisms (so $D_4$ triality is handled),
and the recognition is checked against $\dim K$ computed from the Satake
diagram.
\item $W_H = \Stab_{W_L}(t_j)$ in the adjoint torus, computed by
orbit--stabilizer with Schreier generators; it need not be a reflection
group.
\end{enumerate}

\begin{computation}\label{comp:symmetric}
For all real forms of types A--D (in small rank), $E_6$--$E_8$, $F_4$ and
$G_2$ the construction passes all checks of
Section~\ref{subsec:checks} and the orbit counts of
Section~\ref{subsec:orbitcounts}. Exceptional isomorphisms
($B_2=C_2$, $D_3=A_3$, triality) give identical data from different
diagrams, and rank-one cases agree with the two-orbit front end. Some
dimensions and Euler characteristics:
\begin{center}
\begin{tabular}{lrr@{\qquad}lrr}
\toprule
form & $\dim$ & $\chi$ & form & $\dim$ & $\chi$\\
\midrule
AIII(2,2) & 8 & 39 & DIII(6) & 30 & 4576\\
AIII(3,3) & 18 & 1180 & $G$ & 8 & 27\\
CI(3) & 12 & 148 & FI & 28 & 4788\\
CI(4) & 20 & 1624 & EIV & 26 & 270\\
CII(2,2) & 16 & 123 & EIII & 32 & 2619\\
BI(3,4) & 12 & 147 & AIII(3,4) & 24 & 8015\\
DI(4,4) & 16 & 747 & DIII(7) & 42 & 55168\\
\bottomrule
\end{tabular}
\end{center}
For $G_2/SO_4$, $\chi = 27 = 12+6+6+3$ can be checked by hand.
\end{computation}

\subsection{Toroidal varieties}\label{subsec:toroidal}
Let $X_w$ be wonderful with $\ZZ\Sigma$ as lattice, so that its valuation
cone is the negative orthant in $N=\operatorname{Hom}(\ZZ\Sigma,\ZZ)$ and
the orbit $O_J$ corresponds to the face $F_J$. A smooth fan subdividing the
orthant gives a smooth complete toroidal $X\to X_w$. Its fixed points are
the pairs $(x,\tau)$ with $x\in O_J^T$ and $\tau\subseteq F_J$ a cone of
dimension $\dim F_J$, and near $x_J$ the morphism is the toric modification
of the normal slice, on which $T$ acts through $-\gamma\mapsto\chi_\gamma$.
So the tangent weights at $(w x_J,\tau)$ are $w(\Phi\smallsetminus\Phi_H)$ and
$w(\pr_J(\tau^\vee))$, with $\tau^\vee$ the dual basis. Star subdivisions agree
exactly with the blow-up front end.

\subsection{Orbit counts}\label{subsec:orbitcounts}
Brion and Peyre \cite{BrionPeyre} express the virtual Poincar\'e polynomial
of $G/H$, for $H$ containing a maximal torus, through the Molien series
$F_H$ of $N_H(T)/T$ acting on $\operatorname{Lie}T$:
$|G/H|(q) = q^{\dim G/H}F_H(1/q)/F_G(1/q)$. For each orbit $O_I$ with fixed
points the package compares $|G/P_{S_I}|(q)\cdot|L_{S_I}/H_{L,I}|(q)$ with the
count obtained from the BB cells of the orbit closures by M\"obius
inversion. They agree in all cases tested, including those beyond (R).

% ===========================================================================
\section{Certificates beyond condition (R)}\label{sec:certificates}

For symmetric varieties, $V\ne0$ only for satellites with a Hermitian
factor, and in the cases below $V = \QQ\zeta$ is spanned by the centre of the
Hermitian subgroup. This happens for AIII$(p,q)$ with $p\ne q$, $p\ge2$, for
DIII$(n)$ with $n$ odd, and for EIII; for AIII$(p,p)$ and DIII$(n)$ with $n$
even the component group of $W_H$ negates the centre, and (R) holds. By Proposition~\ref{prop:normal}, the normal weight of
$D_\gamma$ at $x_K$ is $\chi = \pr(-\gamma)+c\,\zeta$ for an unknown rational
$c$. Since all $T$-weights lie in the root lattice, $c$ lies in a coset
$c_0+\ZZ$ that is computed from $\pr(-\gamma)$ and a primitive $\zeta$
(here $c_0=0$). The package uses $c=0$ and certifies it.

\subsection{The point-count certificate}
Let $X^J$ be the closure of an orbit $O_J$ (itself a smooth wonderful
variety), and suppose that all normal weights occurring in $X^J$ are known
except the one of $D_\gamma$ at $x_K$, $K\cup\{\gamma\}\subseteq J$. For a
generic cocharacter $\lambda$ and a fixed point $w x_K$ put
$a_w = \langle w\,\pr(-\gamma),\lambda\rangle$ and
$b_w = \langle w\zeta,\lambda\rangle$. The unknown weight at $wx_K$ pairs to
$a_w+c\,b_w$, so the BB counts $P_\lambda(c)$ of the data with parameter $c$
form a step function of $c$ that jumps only at the thresholds $-a_w/b_w$.

\begin{theorem}\label{thm:certificate}
Let $E(X^J)$ be the E-polynomial of $X^J$, and let $c^*$ be the true value.
\begin{enumerate}
\item If $\lambda$ is generic for the true data (i.e.\ $c^*$ is not a
threshold of $\lambda$), then $P_\lambda(c^*) = E(X^J)$.
\item If $O_J$ has $T$-fixed points, then
\[
E(X^J) = |G/P_{S_J}|(q)\,|L_{S_J}/H_{L,J}|(q) + \sum_{\emptyset\ne S\subseteq J}(-1)^{|S|+1}E(X^{J\smallsetminus S}),
\]
where the first term is given by the Brion--Peyre formula and the others
by the BB counts of closures that do not involve the unknown.
\item In any case $P_\lambda(c^*)$ is palindromic of degree $\dim X^J$.
\end{enumerate}
Consequently, if $\lambda_1,\dots,\lambda_m$ have pairwise disjoint
threshold sets, $m\ge 2$, then $c^*$ lies in
$\bigcap_i\bigl(A_{\lambda_i}\cup T_{\lambda_i}\bigr)$, where $T_\lambda$ is the
set of thresholds and $A_\lambda$ the union of the open intervals between
thresholds on which $P_\lambda$ equals $E(X^J)$ (resp.\ is palindromic).
\end{theorem}
\begin{proof}
(1) With the true weights, the data are the fixed-point data of $X^J$, and
$\lambda$ is generic for them; the BB decomposition of the smooth projective
$X^J$ into affine spaces gives $E(X^J)=\sum_d c_d q^d$.
(2) The orbits of $X^J$ stratify it into the open orbit $O_J$ and the
boundary $\bigcup_{\gamma\in J}X^{J\smallsetminus\{\gamma\}}$, with
$X^{J\smallsetminus S}\cap X^{J\smallsetminus S'}=X^{J\smallsetminus(S\cup S')}$;
the E-polynomial is additive, which gives the inclusion--exclusion. The
orbit $O_J$ fibres Zariski-locally trivially over $G/P_{S_J}$ with fibre
the satellite, and $E(L/H_L)$ is the Brion--Peyre count. The orbit $O_K$
lies in $X^{J\smallsetminus S}$ only if $S\subseteq\{\gamma\}$, and
$X^{K}$ has no normal direction $\gamma$, so no term with $S\ne\emptyset$
involves the unknown.
(3) Poincar\'e duality. For the last statement: every $c$ that is not a
threshold of some $\lambda_i$ with $P_{\lambda_i}(c)\ne E(X^J)$ is excluded by
(1); since the threshold sets are disjoint, $c^*$ is a threshold of at most
one $\lambda_i$, so it is decided by the others.
\end{proof}

Intersecting with $c_0+\ZZ$ gives a finite list of candidates when every
admissible interval is bounded; if only $c=0$ remains, $c=0$ is proved,
given the fixed-point data of the known orbits. In practice the admissible
set at each $\lambda$ is a single interval around $0$ of half-width below
$0.1$.

Two unknowns on different orbits are handled jointly: every fixed point
sees at most one of them, so $P_\lambda = P_0+F_1(c_1)+F_2(c_2)$, and the
admissible boxes are found by matching the values of $F_2$ against
$E-P_0-F_1$ in a hash table.

\subsection{Induction over orbit closures}
The unknowns $(K,\gamma)$ are decided one at a time, each on the smallest
closure $X^J\supseteq X^{K\cup\{\gamma\}}$ whose other normal weights are
proved or already certified, preferring closures whose open orbit has fixed
points. If the palindromic condition leaves $c$ unbounded, the next larger
closure is tried; if no single unknown can be decided, pairs are decided
jointly.

\begin{computation}\label{comp:certified}
The certificate leaves exactly $c=0$ for every unknown normal weight of the
complete symmetric varieties AIII(2,3), AIII(2,4), AIII(2,5), AIII(3,4),
AIII(3,5), AIII(3,6), AIII(4,5), DIII(5), DIII(7) and EIII.
For AIII$(3,4)$ the orbit $O_{13}$ has no fixed points, and the unknown
$(O_3,\gamma_0)$ is decided on $X^{\{0,2\}}$ by palindromy alone. For
DIII$(7)$ palindromy leaves $c$ unbounded there, and $(O_3,\gamma_0)$,
$(O_{23},\gamma_0)$ are decided jointly on $X$. For AIII$(4,5)$ all six
unknowns are decided one at a time.
\end{computation}

\begin{remark}
A first certificate used the Atiyah--Bott--Berline--Vergne identity
$\sum_p 1/e_p(\lambda)=0$ \cite{AtiyahBott84,BerlineVergne} as a rational
function of $c$: an explicit bound excludes large $|c|$, and all integers
below it are screened modulo a large prime. It certifies the cases
AIII$(2,q)$, DIII$(5)$ and EIII, but for AIII$(3,4)$ and DIII$(7)$ the
coefficients of the expansion in $1/c$ vanish identically and the moments
$\int c_1^T(TX)^k$ give bounds beyond reach. This degeneracy motivated
Theorem~\ref{thm:certificate}.
\end{remark}

\begin{question}
Is $c=0$ for every Hermitian satellite, i.e.\ is the normal weight always the
$W_L$-average for complete symmetric varieties? A sketch: the Satake
involution fixes every spherical root and swaps the two end colours of a
Hermitian block, so the colour coefficients $c(D,\gamma)$ are invariant, and
the class measuring $c$ vanishes. Making this rigorous requires the sign
of $\operatorname{div} f_\gamma$, the reducedness of $D\cap F$, and that
colours not moved by $S_I$ miss the localization.
\end{question}

% ===========================================================================
\section{Cell counts without listing fixed points}\label{sec:large}

For a generic $\lambda$ the cell dimension of $w x_O$ is the number of
$\lambda$-positive weights among $w(\Phi\smallsetminus\Phi_H)$ and $w(N_O)$.
The coset search of Section~\ref{sec:spherical} enumerates the minimal
representatives level by level, keeping only one length level of inverse
permutations of the roots in memory. When $W_H$ is not generated by
reflections, the counts over cosets of the reflection part are divided by
$|W_H:W_{\mathrm{refl}}|$, and the division must be exact. Two
pseudo-random cocharacters must give the same counts, and the counts must
be palindromic.

\begin{computation}\label{comp:large}
The complete symmetric varieties of the following exceptional types have
the stated dimension and Euler characteristic; the counts are
$\lambda$-independent and palindromic (times on one core in pure Python):
\begin{center}
\begin{tabular}{lrrr}
\toprule
form & $\dim$ & $\chi$ & time\\
\midrule
EII & 40 & 110\,916 & 4\,s\\
EI & 42 & 370\,170 & 14\,s\\
EVII & 54 & 23\,464 & 7\,s\\
EVI & 64 & 758\,079 & 48\,s\\
EV & 70 & 28\,373\,976 & 30\,min\\
EIX & 112 & 7\,445\,880 & 13\,min\\
\bottomrule
\end{tabular}
\end{center}
For example, the cell counts of EVII are 1, 3, 7, 12, 19, 29, 44, 63, 87,
117, 155, 198, 248, 303, 366, 432, 503, 575, 652, 724, 795, 858, 920, 970,
1013, 1041, 1063, 1068, \dots, symmetric around the middle term 1068. EVIII is out of reach this way:
its closed orbit alone has $|W(E_8)| = 696\,729\,600$ fixed points.
\end{computation}

% ===========================================================================
\section{Equivariant cohomology beyond GKM}\label{sec:cohomology}

For GKM varieties, $H_T^*(X;\QQ)$ is the ring of tuples of polynomials
$(f_p)_{p\in X^T}$ with $f_p\equiv f_q\bmod\chi$ along every $T$-curve
\cite{GKM}. Complete quadrics are not GKM: at a fixed point the weights
$\alpha$ and $2\alpha$ can both occur, and there are infinitely many
$T$-curves.

\subsection{Brion's description}
For a smooth complete spherical variety, Brion \cite[Thm.~6 and proof of
Thm.~9]{Brion98} cuts out the image of $H_T^*(X)\to H_T^*(X^T)$ by
conditions attached to the connected components $Y$ of $X^{\ker\chi}$ for
primitive characters $\chi$. Such a $Y$ is a point, a $\PP^1$ (condition
$f_y\equiv f_z\bmod\chi$), or a surface, namely $\PP(\mathfrak{sl}_2)$ or a
rational ruled surface; for a surface all $f_p$, $p\in Y^T$, are congruent
modulo $\chi$ and $\sum_{p\in Y^T} f_p/e_p$ is a polynomial, where $e_p$ is
the product of the two weights of $T_pY$.

\begin{lemma}\label{lem:components}
The components can be read off the fixed-point data: points of one
component of $X^{\ker\chi}$ have the same weights modulo $\ZZ\chi$, and the
weights along $\chi$ form one of the patterns
$\{k,-k\}$ (curve), $\{a,-a\},\{a,2a\},\{-a,-2a\}$ (plane), or
$\{\pm a,\pm b\}$ (ruled surface).
\end{lemma}
The package raises when the grouping is ambiguous; it never was in the
examples. On flag and toric varieties the inferred curves are the GKM
edges; complete conics have $12$ curves and $6$ planes, complete quadrics
in $\PP^3$ have $153$ curves and $48$ planes. In low degrees the ring cut
out by these conditions has the dimension of a free module with the BB
Poincar\'e series, while dropping the second-order condition gives too many
classes.

\subsection{Characteristic numbers}
A degree-one class is determined by its values on the closed orbit:
propagation along the components determines the other values, and all
conditions are checked at the end. A $G$-linearized line bundle with weight
$\mu$ at the base point $z$ of the closed orbit has value $w(\mu)$ at $wz$.
For complete quadrics in $\PP^{n-1}$, the colours $\mu_k$ are the pullbacks of
$\mathcal O(1)$ from $\PP(\mathrm{Sym}^2\Lambda^kV)$, of weight $-2\omega_k$.
Monomials in these classes are integrated by the localization formula
\cite{AtiyahBott84,BerlineVergne} in exact arithmetic; the result must be an
integer independent of the cocharacter. The class of the condition to be
tangent to a given quadric is $2(\mu_1+\dots+\mu_{n-1})$.

\begin{computation}\label{comp:characteristic}
The package reproduces the following numbers for complete conics,
complete quadric surfaces and complete quadrics in $\PP^4$ and $\PP^5$;
$\mu,\nu,\rho$ and $\mu_1,\dots,\mu_{n-1}$ are the colours.

\smallskip
\begin{center}
\small
\begin{tabular}{p{0.2\textwidth}p{0.33\textwidth}p{0.36\textwidth}}
\toprule
$X$ & number & value\\
\midrule
conics & $\mu^a\nu^{5-a}$, $a=5,\dots,0$ & 1, 2, 4, 4, 2, 1\\
 & $(2\mu+2\nu)^5$ (Chasles \cite{Chasles}) & 3264\\
quadric surfaces & $\mu^a\nu^{9-a}$, $a=9,\dots,0$ \cite{Schubert} & 1, 2, 4, 8, 16, 32, 56, 80, 92, 92\\
 & $\rho^9$ & 1\\
 & $(2\mu+2\nu+2\rho)^9$ & 666\,841\,088\\
 & $\mu^3\nu^3\rho^3$, $\mu^2\nu^5\rho^2$ \cite{BFS} & 104, 128\\
quadrics in $\PP^4$ & $(2\sum\mu_k)^{14}$ \cite{Sturmfels3264} & 48\,942\,189\,946\,470\,400\\
 & $\mu_2^{14}$ & 7703\\
quadrics in $\PP^5$ & $(2\sum\mu_k)^{20}$ & see below\\
 & $\mu_k^{20}$, $k=1,\dots,5$ & 1, 803128, 61520094, 803128, 1\\
\bottomrule
\end{tabular}
\end{center}
\smallskip

\noindent It also reproduces the maximum likelihood degrees
$\varphi(n,d)=\int\mu_1^{\binom{n+1}2-d}\mu_{n-1}^{d-1}$ of generic linear
concentration models \cite[Prop.~3.5]{MMMSV}, as tabulated in
\cite[\S2.2]{SturmfelsUhler}: $\varphi(4,d)$ = 1, 3, 9, 17, 21, 21, 17, 9, 3, 1;
$\varphi(5,d)$ = 1, 4, 16, 44, 86, 137, 188, 212, \dots; and
$\varphi(6,d)$ = 1, 5, 25, 90, 240, 528, 1016 for $d\le7$. For complete
quadrics in $\PP^5$ ($3690$ fixed points, $28\,440$ components) the number of
quadrics tangent to $20$ general quadrics comes out as
\[
\int_X \bigl(2(\mu_1+\dots+\mu_5)\bigr)^{20} = 1\,810\,718\,299\,257\,984\,458\,113\,941\,504 .
\]
\end{computation}

\begin{remark}
The list in \cite[Question~5]{Sturmfels3264} gives
\[641\,211\,464\,734\,373\,953\,791\,690\,014\,720\]
for this number. A
second computation by a different method, not yet part of the package,
gives our value, and so does the growth of the sequence: the decimal logarithms of the four numbers
are $3.51$, $8.82$, $16.69$ and $27.26$ (ours), with differences $5.31$,
$7.87$, $10.57$, whereas the listed value would give $29.81$. We have not
resolved the discrepancy; a third computation is planned.
\end{remark}

\subsection{Cell classes}
Call a class $\tau_p\in H_T^*(X)$ \emph{canonical} (after Goldin--Tolman
\cite{GoldinTolman}) if $\tau_p(p)$ is the product $e^-_p$ of the
$\lambda$-negative weights at $p$ and $\tau_p(q)=0$ for all $q\ne p$ with
$\codim q\le\codim p$, where $\codim p$ is the number of negative weights.

\begin{proposition}\label{prop:canonical}
A canonical class for $p$ is unique if it exists. It exists if every
fixed point $q\ne p$ in the closure of the plus-cell $C_p$ satisfies
$\dim C_q<\dim C_p$; then it is the equivariant class of
$\overline{C_p}$. Conversely, if the canonical class of $p$ does not exist,
then $\overline{C_p}$ contains a fixed point $q\ne p$ with $\dim C_q\ge \dim C_p$.
\end{proposition}
\begin{proof}
The class $[\overline{C_p}]$ restricts to $e^-_p$ at $p$, since $C_p$ is
smooth there with normal space the negative weight space, and to zero at
fixed points outside $\overline{C_p}$. If $\overline{C_p}$ contains no fixed
point $q\ne p$ with $\codim q\le\codim p$, this class is canonical; this
proves existence and, by contraposition, the last statement. Uniqueness:
the difference $\delta$ of two canonical classes vanishes at $p$ and at all
$q$ with $\codim q\le\codim p$. Order the fixed points by a Morse function of
an ample class; if $\delta\ne0$, let $q$ be minimal with $\delta(q)\ne 0$. By
the Atiyah--Bott lemma \cite{AtiyahBott84}, applied to the filtration by
closed unions of plus-cells (whose open parts have affine pavings and are
therefore equivariantly formal), $\delta(q)$ is divisible by $e^-_q$, which has degree $\codim q>\codim p=\deg\delta$,
a contradiction.
\end{proof}

The package constructs canonical classes point by point in a Morse order.
At $q$, each component $Y\ni q$ on which $q$ has $m\in\{1,2\}$ negative weights
gives a congruence $\tau(q)\equiv L_Y\bmod\chi_Y^m$, with $L_Y$ the value at an
earlier point ($m=1$) or obtained from the sum condition ($m=2$); the
moduli multiply to $e^-_q$, and $\tau(q)$ is found by Newton interpolation
using only restriction to hyperplanes and exact division. On flag varieties
the result equals the GKM flow-up classes, and on the non-GKM examples
CII$(1,2)$ and BI$(1,4)$ it equals a global linear solve.

\begin{computation}\label{comp:nocanonical}
For complete conics the line arrangement of the tangent weights has $12$
chambers. In each chamber exactly two of the twelve fixed points have no
canonical class. Hence, by Proposition~\ref{prop:canonical}, for every
generic cocharacter two plus-cell closures contain fixed points of cells
of at least their own dimension: the BB decomposition of complete conics is
never a stratification by cell closures. Canonical classes also fail to exist
for some points of AI(4), CI(2), CI(3), $G_2$, AIII(2,2), DI(2,4) and even of
the GKM variety AII(3) for the default cocharacter, and they exist for all
points of AIII(1,2), BI(1,4) and CII(1,2).
\end{computation}

\subsection{The divisor subalgebra}
If $A = \bigoplus_{d=0}^n A^d$ is a graded Poincar\'e duality algebra over
$\QQ$ generated by $A^1$, then $A\cong\QQ[x_1,\dots,x_k]/\Ann(V)$ with the
volume polynomial $V(x) = \int(\sum_i x_iD_i)^n/n!$, acting by
differentiation (Macaulay's inverse systems \cite{Macaulay}). For
$D_i\in H^2(X;\QQ)$ the same construction yields the subalgebra generated by
the $D_i$, whose Hilbert function is computed from the intersection numbers
$\int D^a$: a form $f$ of degree $d$ vanishes in $H^*(X;\QQ)$ if and only if
$\int fg=0$ for all $g$ of degree $n-d$.

\begin{computation}\label{comp:divisors}
\begin{enumerate}
\item Complete conics: $H^*(X;\QQ)=\QQ[\mu,\nu]/(r_3,r_4)$ with one cubic
relation $r_3 = 2\nu^3-3\mu\nu^2+3\mu^2\nu-2\mu^3$ and one further quartic
relation.
\item Complete quadric surfaces: $H^*(X;\QQ)$ is generated by $\mu,\nu,\rho$,
with minimal relations in degrees $4,4,5,5,6$.
\end{enumerate}
\end{computation}

\begin{proposition}\label{prop:notgenerated}
The cohomology of complete quadrics in $\PP^4$ is not generated by divisor
classes.
\end{proposition}
\begin{proof}
The Picard rank is $4$, so $H^2$ has dimension $4$ and the degree-$6$ part of
the subalgebra generated by $H^2$ has dimension at most $\binom{6}{3}=20$,
whereas $b_6=21$ (the cell counts are $1,4,10,21,36,53,65,70,\dots$).
\end{proof}
The computation shows more: the subalgebra generated by the colours has
Hilbert function $1,4,10,20,35,52,65,70,\dots$. The missing classes are the
subject of \cite{DGMP}.

% ===========================================================================
\section{Real points}\label{sec:real}

\subsection{The question}
In the stable motivic homotopy category, the component of the attaching map
from a cell of dimension $d+1$ to a cell of dimension $d$ of a motivic cell
structure lies in $\pi_{1,1}=W\cdot\eta$ \cite{MorelVoevodsky,DuggerIsaksen}.
For smooth projective varieties over $\ZZ$ with a BB decomposition we ask
whether it comes from $W(\ZZ)\cdot\eta\cong\ZZ\cdot\eta$ (hypothesis H1 of the
project). Under real realization $\eta$ has degree $\pm2$, so the incidence
numbers of the real cell complex of $X(\RR)$ would be twice these
coefficients; under complex realization $\eta$ is detected by $Sq^2$.

\subsection{The rule}
For a fixed point $p$ let $\sigma(p)$ be the sum of the $\lambda$-positive
weights at $p$. For an invariant curve from $x$ to $y$ with
$d_\lambda(x)=d_\lambda(y)+1$ and weight $\varphi$ at $x$, write
$\sigma(x)-\sigma(y)=m\,\varphi$.

\begin{conjecture}[H1${}'$]\label{conj:h1}
Let $X$ be smooth projective over $\ZZ$ with a split torus action, and let
$\lambda$ be a generic cocharacter whose decomposition is \emph{graded}:
every invariant curve goes down in cell dimension. Then the
$\eta$-component of the attaching map from $x$ to $y$ is $\pm\eta$ if $x$
and $y$ are joined by an invariant curve with $m$ even, and $0$ otherwise.
\end{conjecture}

For $G/P$ with $\lambda$ dominant regular, $\sigma$ is Kocherlakota's
$\sigma$ and the real statement is Kocherlakota's theorem \cite{Kocherlakota};
the package checks the identification for every parabolic of
$A_2,A_3,B_2,B_3,C_3,G_2,D_4$. For toric varieties the predicted
incidences, completed by the unique signs (up to reorientation) with
$\partial^2=0$, give the rational Betti numbers of $X(\RR)$ computed
independently by the formula of Choi--Park and Suciu--Trevisan
\cite{ChoiPark,SuciuTrevisan}. In the complex realization, the parity rule
agrees with $(y^*)^2\bmod2$ computed in the GKM ring for all examples tested.

\subsection{Exact incidences for toric varieties}
The real points of a smooth complete toric variety with polytope $P$ are
$X(\RR)=P\times(\ZZ/2)^n/\!\sim$ \cite{DavisJanuszkiewicz}. A face of $P$,
i.e.\ a cone $\tau$ of the fan, contributes one cell of dimension
$n-|\tau|$ for each class of $(\ZZ/2)^n$ modulo the span $\Lambda_\tau$ of
the rays of $\tau$ modulo $2$. Orienting each cell by its face, the
attaching maps are the identity on $P$ and
\[
\partial(\tau,e)=\sum_{r}(-1)^{\,n-|\tau|+\#\{i\in\tau:\,i<r\}}\,(\tau\cup\{r\},e),
\]
summed over the rays $r$ with $\tau\cup\{r\}$ a cone. The real BB cell of a
vertex $v$ is the union of the cells of the faces whose $\lambda$-minimal
vertex is $v$; in the ray coordinates of $v$ it is the set $\{+,-,0\}^{d_v}$
of sign vectors.

\begin{theorem}\label{thm:morse}
Match a cell of the BB cell of $v$ whose first coordinate (in a fixed order)
that is not $+$ equals $0$ with the cell obtained by replacing that
coordinate by $-$. This is an acyclic matching of the Davis--Januszkiewicz
complex whose critical cells are the positive orthants, one for each fixed
point, of dimension $d_\lambda(v)$. Consequently the algebraic Morse complex
\cite{Forman,Skoldberg,Kozlov} is a chain complex of free abelian groups with
one generator of degree $d_\lambda(v)$ for each fixed point, whose homology is
$H_*(X(\RR);\ZZ)$.
\end{theorem}
\begin{proof}
Within one BB cell the matching is the standard lexicographic collapse of
$\{+,-,0\}^d$ onto $(+,\dots,+)$, and matched cells are faces of each other
of codimension one with incidence $\pm1$. Gradient paths leave the BB cell
of $v$ only towards faces not containing $v$, whose minimal vertex has a
larger value of $\langle\cdot,\lambda\rangle$; so a closed path would have to
stay in one BB cell, where the matching is acyclic. The last statement is
the main theorem of algebraic Morse theory.
\end{proof}

If the decomposition is graded, every vertex of $F_v$ other than $v$ is
reached from $v$ by a monotone edge path, along which the cell dimension
drops; hence $\overline{C_v(\RR)}\smallsetminus C_v(\RR)$ lies in cells of
smaller dimension, the real BB cells form a CW structure, and the
incidences can be compared entry by entry.

\begin{observation}\label{obs:entrywise}
On random smooth complete toric varieties with graded decompositions, the
magnitude of every entry of the Morse complex equals the prediction of
Conjecture~\ref{conj:h1} ($2$ if and only if the cells are joined by a curve
with $m$ even): $14$ decompositions in dimension $2$, $42$ in dimension $3$
and $40$ in dimension $4$, with $9$, $70$ and $151$ nonzero incidences, and no
mismatch.
\end{observation}

\subsection{Beyond GKM}
For non-GKM varieties we apply the rule to the invariant curves provided by
the components of Lemma~\ref{lem:components}: every $\PP^1$ component, the
two lines of weight $\pm a$ in each plane $\PP(\mathfrak{sl}_2)$ (the conics of
weight $2a$ join cells two apart), and the four boundary curves of each
ruled surface. The parity $m$ is computed after pairing the normal weights at
the two ends of the curve; when normal weights are congruent modulo
$\varphi$ the pairing is a choice, but flip-free pairings all give the same
$m$.

\begin{computation}\label{comp:beyondgkm}
For each of the following varieties and $12$ random generic cocharacters,
the decomposition is graded along these curves, the sign completion is
unique, and the predicted rational Betti numbers of $X(\RR)$ equal the known
ones:
\begin{center}
\small
\begin{tabular}{lll}
\toprule
variety & torus & known $X(\RR)$\\
\midrule
$\PP^q\times\check\PP^q$, $q=2,3,4$ & $PGL_{q+1}$ & $\RP^q\times\RP^q$\\
$\Gr(2,6)$, $\Gr(2,8)$ & $Sp_6$, $Sp_8$ & \cite{CasianKodama}\\
$E_6/P_1$ & $F_4$ & \cite{Kocherlakota}, $E_6$-torus\\
\bottomrule
\end{tabular}
\end{center}
In each case $X(\RR)$ is the same real variety as for a larger torus with
GKM data, where the answer is known. In the notation of symmetric varieties
these are AIII$(1,q)$, CII$(1,2)$, CII$(1,3)$ and FII.
\end{computation}

\subsection{Non-graded decompositions}
Every rank-two complete symmetric variety we tested (AI(3), AI(4), AII(3),
AIII(2,2), AIII(2,3), BI(2,3), CI(2), CI(3), CII(2,2), DI(2,4), $G_2$) is
non-graded along the invariant curves for all $576$ random cocharacters
tried, and already the toric surface $dP_6$ has no graded cocharacter. With
rational coefficients a chain complex on the BB cells still exists, but
only up to filtered change of basis.

\begin{observation}[surfaces]\label{obs:deflection}
Let $a\to b$ be a curve between cells of equal dimension; exactly one normal
direction $\psi$ changes sign along it, and put
$m=(\sigma(a)-\sigma(b)+\psi_b)/\varphi$. Add to the direct terms of
Conjecture~\ref{conj:h1} a term $\pm2$ in $\langle\partial c,b\rangle$ for
every curve $c\to a$ with $d(c)=d(a)+1$, whenever $m$ is odd. Heuristically,
each of the two real arcs from $b$ to $a$ is deflected into one ascending arc
of $a$. For $dP_6$, four other non-graded fans and $144$ decompositions of
$30$ random toric surfaces, every sign completion gives the Choi--Park Betti
numbers. In dimension $3$ the rule fails, and the exact complexes of
Theorem~\ref{thm:morse} have nonzero entries between cells that share no
neighbour.
\end{observation}

\begin{proposition}\label{prop:realconics}
For the variety $X$ of complete conics,
$H^*(X(\RR);\QQ)\cong H^*(\RP^5;\QQ)$, so the rational Betti numbers of
$X(\RR)$ are $1,0,0,0,0,1$.
\end{proposition}
\begin{proof}
$X=\Bl_{v_2(\PP^2)}\PP^5$, so $X(\RR)$ is the real blow-up of $\RP^5$ along
the Veronese $\RP^2$, with exceptional divisor $E$ an $\RP^2$-bundle over
$\RP^2$. The blow-up square gives a Mayer--Vietoris sequence
\[\cdots\to H^k(\RP^5)\to H^k(X(\RR))\oplus H^k(\RP^2)\to H^k(E)\to\cdots.\]
Since $\RP^2$ is $\QQ$-acyclic, $E\to\RP^2$ induces an isomorphism in rational
cohomology, and the five lemma gives the claim.
\end{proof}

\begin{observation}\label{obs:conics}
For complete conics, the prediction of Conjecture~\ref{conj:h1} along the
invariant curves admits no sign completion in any chamber. In each of $21$
chambers tested, removing exactly one term gives a complex whose rational
homology is that of Proposition~\ref{prop:realconics}, and this term is unique:
it runs from the upper end of the curve between cells of equal dimension
with $m$ even to the lower end of the one with $m$ odd (there are always
exactly two such curves). Neither the rule of
Observation~\ref{obs:deflection} nor a sum over paths of such curves
explains this; the planes $\PP(\mathfrak{sl}_2)$ seem to need their own local
model.
\end{observation}

% ===========================================================================
\section{Software}\label{sec:software}

\bbcells{} is free software under the GNU General Public License, version~3.
It is written in Python ($\ge3.10$) with the standard library only;
all arithmetic is exact (\code{int}, \code{Fraction}, and modular arithmetic
for probabilistic identity tests). The core type is the fixed-point data;
front ends produce them from fans, root data, Satake diagrams, spherical
data or constructions (products, restriction to subtori, blow-ups), and the
invariants, the equivariant cohomology and the real complexes are computed
from them. The command line covers the front ends and invariants, for
example
\begin{verbatim}
python3 -m bbcells symmetric AIII 2 3
python3 -m bbcells symmetric EVI --counts-only
python3 -m bbcells real A3
\end{verbatim}
The test suite has more than $300$ tests, including doctests, and runs in
about a minute; it contains the oracle tests of
Table~\ref{tab:frontends} and every computation of this paper that takes
less than a few seconds. The larger computations are reproduced by
\code{paper/reproduce.py}.

% ===========================================================================
\section{Open problems}\label{sec:open}

\begin{enumerate}
\item \emph{General Luna data.} Wonderful varieties whose satellites are not
symmetric, starting with Wasserman's rank-two list \cite{Wasserman} and the
primitive spherical systems of Bravi--Pezzini \cite{BravPez,BravPezPrimitive}.
\item \emph{The Hermitian case in general} (the question in
Section~\ref{sec:certificates}), and certificates for two unknowns on one
orbit.
\item \emph{The number of quadrics in $\PP^5$ tangent to $20$ general quadrics}
(Remark after Computation~\ref{comp:characteristic}).
\item \emph{An integral basis} of $H^*(X;\ZZ)$ adapted to the cells when
canonical classes do not exist, as for complete conics; and the generators
beyond the divisors for complete quadrics in $\PP^4$.
\item \emph{Toroidal varieties with $\Lambda\supsetneq\ZZ\Sigma$}, which are
finite covers of toroidal varieties over wonderful models.
\item \emph{Non-graded real incidences} in dimension $\ge3$
(Observations~\ref{obs:deflection} and~\ref{obs:conics}), and a proof of
Conjecture~\ref{conj:h1} in $SH(\ZZ)$, for instance through the Thom-space cell
structures of \cite{VoelkelSpherical} restricted to invariant curves.
\item \emph{EVIII}, via the orbit decomposition rather than the fixed points.
\end{enumerate}

\subsection*{Acknowledgements}
\emph{[To be written by the author; see the comment in the source.]}
% To be decided by the author (PLAN.md 4.4): acknowledgement of tool use.
% Parts of the software and a first draft of this text were written with
% the help of an AI coding assistant (Claude Code); all statements are to be
% checked by the author before submission.

\bibliographystyle{amsplain}
\bibliography{references}

\end{document}
